Keberadaan sebuah ruang topologis $E$ dengan sifat-sifat tersebut mengikuti dari
[[Topologischer Raum/Verklebungsdatum/Existenz/Fakt|fakta tentang keberadaan ruang hasil perekatan]]
\zusatzklammer {himpunan-himpunan terbuka yang harus direkatkan adalah

\mavergleichskettek
{\vergleichskettek
{W_{ij}
}
{ \defeq }{ E_i \vert_{U_i \cap U_j }
}
{  }{
}
{    }{
}
{   }{
}
}
{}{}{}} {} {}
dan keberadaan pemetaan kontinu ke $X$ mengikuti dari
[[Topologischer Raum/Verklebungsdatum/Stetige Abbildung/Fakt|fakta tentang pemetaan kontinu dari ruang hasil perekatan]].
Pada setiap serat $E_x$ terdapat struktur ruang vektor yang terdefinisi dengan baik, yang berasal dari $E_i$ untuk sebarang lingkungan terbuka

\mavergleichskette
{\vergleichskette
{ x
}
{ \in }{ U_i
}
{  }{
}
{    }{
}
{   }{
}
}
{}{}{}.
Ketaktergantungan terhadap pilihan tersebut didasarkan pada kenyataan bahwa untuk

\mavergleichskette
{\vergleichskette
{ x
}
{ \in }{ U_i \cap U_j
}
{  }{
}
{    }{
}
{   }{
}
}
{}{}{}
menurut asumsi terdapat sebuah isomorfisme bundel vektor
\maabbdisp {\varphi_{ji}} {E_i \vert_{U_i \cap U_j } } { E_j \vert_{U_i \cap U_j }
} {}
yang menginduksi sebuah
\definitionsverweis {isomorfisme ruang vektor}{}{}
\maabbdisp {} { \left(  E_i  \right)_x } { \left(  E_j  \right)_x
} {}.

 [[Kategorie:Latexseite]]
